\chapter{Análisis}

Una hoja en blanco es un mundo de posibilidades para algunos, si embargo para mi
es una situación intimidante, la recomendación es adquirir realimentación tan 
rápido como sea posible y la fuente más fiable es aquel que tiene el problema, 
para mi fortuna tengo la asistencia del laboratorio del \emph{Laboratorio de 
Investigación y Desarrollo de Aplicaciones Interactivas para la
Neuro-Rehabilitación} del Instituto de Fisiología Celular de la UNAM, quienes
son expertos y tienen gran experiencia tratando y rehabilitando a este tipo 
de pacientes, gran parte del entendimiento y realimentación es gracias a ellos.

El problema habla por si solo, pero para entenderlo se requiere hablar su idioma,
se le conoce como lenguaje del dominio y la forma de entenderlo es hablando 
con las personas adecuadas, no se trata de socializar si no de entender mejor 
las necesidades y requerimientos.

Para esta primera etapa sabemos que los pacientes son personas que han perdido 
la movilidad de los brazos, las terapias de recuperación consisten en ejercicios
de movimiento repetitivos, inducir movimientos como desplazamientos  y rotaciones
en lo brazos de los pacientes provocan una mejoría motriz sostenida.   

Para mi caso lo que se me requiere es crear un dispositivo que pueda asistir en
desplazamientos sobre el plano horizontal.

Hay configuraciones de dispositivos que pudieran realizar esta tarea, por lo que
considero que tomar una de ellas y adaptarla a nuestras necesidades es la opción
correcta. 

\section{Selección de robot}
Dentro de los dispositivos disponibles la elección era obvia, un robot, ¿pero  de
que tipo?, me decante por tres:
\begin{description}
	\item[Robot antropomorfico] sin duda es el robot más versátil que existe,
		es ampliamente utilizado en la industria manufacturera sobre todo
		en la automotriz, sin embargo tiene una gran desventaja para 
		nuestro propósito, es muy costoso de implementar, principalmente
		por el tipo de actuadores que requiere, esta configuración produce 
		un torque muy grande en su base, lo que deriva en una base de 
		acero y actuadores hidráulicos o neumáticos.
	\item[Robot prismático] Este robot es excelente para movimientos precisos 
		en el plano, es fácil de implementar y soporta cargas muy grandes,
		en un principio esta era una opción muy atractiva. Tiene un par 
		de desventajas, es relativamente caro y es menos mantenible que 
		la tercer opción.
	\item[Robot móvil] Es la opción que satisface de mejor manera el 
		requerimiento, se puede desplazar sobre trayectorias largas, 
		rectas o circulares, es rápido, ligero y de los tres es el más
		barato, no tiene la precisión del robot prismático pero no la 
		requerimos, pocas partes mecánicas y ampliamente documentado. 
		En contraste es más difícil de programar pues requiere más 
		sensores que los anteriores.
\end{description}

\section{Requerimientos}
El objetivo del análisis es convertir el problema del lenguaje del dominio a un
problema en el lenguaje de la ingeniería, como ya se menciono el proyecto es 
iterativo y los requerimientos cambiaran constantemente con el propósito de 
ser exhaustivo en su cumplimiento, los requerimientos para esta primera iteración
se presentan en el capitulo primera iteración.
